\title{Group Sequence Policy Optimization}

\begin{abstract}
We present Group Sequence Policy Optimization (GSPO), a reinforcement learning method designed for robust, efficient, and high-performing large language model training. Departing from prior approaches that utilize token-level importance ratios, GSPO instead constructs the importance ratio on the likelihood of entire sequences, applying sequence-level strategies for clipping, rewards, and optimization. Through extensive evaluation, we show that GSPO surpasses the GRPO algorithm in both training efficiency and model performance, provides notable stabilization during Mixture-of-Experts (MoE) RL training, and can streamline the architecture of RL systems. The advantages of GSPO have played a key role in driving the recent advances observed in the Qwen3 models.
\end{abstract}

\section{Introduction}
\label{sec:intro}

Reinforcement learning (RL) has become a central approach for advancing the capabilities of language models \citep{o1,dpsk-r1,qwq32b,qwen3}. By applying RL at scale, these models acquire the proficiency to address complex tasks, including competition-grade mathematics and programming, engaging in more extended and intricate reasoning steps.

For RL methods to effectively leverage increased computational resources, ensuring consistent and resilient training processes is essential. Yet, leading RL approaches such as GRPO \citep{grpo} are plagued by notable instability when applied to training very large language models, frequently culminating in catastrophic and unrecoverable model failures \citep{qwen3, minimax-m1}. Such instability represents a significant obstacle to advancing the capabilities of language models via extended RL-based training.

In this study, we demonstrate that the core instability present in GRPO arises from a foundational misutilization and breakdown of importance sampling weights within the architecture of the algorithm. This flaw leads to the injection of substantial high-variance noise into the training process, with this effect compounding as response lengthens, and the situation is exacerbated by the use of clipping, culminating in model failure.

To overcome these fundamental challenges, we introduce \textbf{Group Sequence Policy Optimization (GSPO)}, a reinforcement learning framework specifically designed for optimizing large language models. GSPO’s principal contribution is its theoretically motivated formulation of the importance ratio, which relies on sequence likelihood \citep{click} and adheres to the standard concept of importance sampling. Furthermore, GSPO derives normalized rewards by assessing the advantages among various responses to a given query, thereby maintaining coherence between rewards at the sequence level and the optimization process.

Our experimental results reveal that GSPO consistently outperforms GRPO with respect to training stability, efficiency, and overall performance. Notably, GSPO naturally addresses the stability issues commonly encountered when applying reinforcement learning to large-scale Mixture-of-Experts (MoE) models, removing the reliance on elaborate stabilization techniques and potentially streamlining RL system design. These strengths of GSPO have played a pivotal role in driving the notable performance enhancements observed in the recent Qwen3 models. We anticipate that GSPO will serve as a reliable and extensible basis for ongoing progress in large-scale RL-based language model training.

\section{Preliminaries}

\paragraph{Notation}
Throughout this work, we represent an autoregressive language model with parameters $\theta$ as a policy, denoted by $\pi_\theta$. The variable $x$ refers to a single query, and $\mathcal{D}$ stands for the set of all queries. For a given query $x$, if the model generates a response $y$, the probability of $y$ under the policy $\pi_\theta$ is expressed as $\pi_\theta (y | x)=\prod_{t=1}^{|y|} \pi_\theta (y_t | x, y_{<t})$, where $|y|$ signifies the token count in $y$. Each query-response pair $(x, y)$ may be evaluated using a verifier $r$, which assigns a reward $r(x, y) \in [0, 1]$.

\paragraph{Proximal Policy Optimization (PPO)} PPO \citep{ppo} updates the policy using data collected from the previous policy $\pi_{\theta_\text{old}}$, employing a clipping strategy to restrict policy changes to a neighborhood around the old policy. More precisely, PPO maximizes the following surrogate objective for policy optimization (here, we exclude the KL regularization component for simplicity, since it is not central to our discussion):

\begin{align}
\mathcal{J}_\text{PPO}(\theta) = \mathbb{E}_{ x \sim \mathcal{D},\, y \sim \pi_{\theta_\text{old}}( \cdot | x) }
\left[ \frac{1}{|y|} \sum_{t=1}^{|y|} 
\min \left( w_{t}(\theta) \widehat{A}_{t},  \, \mathrm{clip} \left( w_{t}(\theta), 1 - {\varepsilon}, 1 + {\varepsilon}\right) \widehat{A}_{t} \right)
\right],
\end{align}

Here, the token $y_t$'s importance ratio is given by $
w_{t}(\theta) = \frac{ \pi_{\theta} (y_{t} | x, y_{<t}) }{ \pi_{\theta_\text{old}} (y_{t} | x,y_{<t})} $. The advantage $\widehat{A}_{t}$ corresponding to $y_t$ is approximated using a separate value model, and $\varepsilon$ denotes the threshold for clipping the importance ratios.

A fundamental practical limitation of PPO stems from its significant dependence on the value model. Typically, the value model matches the policy model in scale, thereby substantially increasing both memory usage and computational demands. Moreover, the performance of the algorithm is closely tied to the accuracy of the value estimation it provides. Obtaining an accurate value model is already a difficult endeavor, and extending this reliability to handle lengthier outputs or more sophisticated tasks exacerbates the problem.

\paragraph{Group Relative Policy Optimization (GRPO)} GRPO \citep{grpo} eliminates the reliance on a value function by directly assessing the relative advantage of individual responses in relation to other responses within the same group for a given query. In particular, GRPO seeks to maximize the following objective:

\begin{align}
\label{equ:grpo}
\mathcal{J}_\text{GRPO}(\theta) = \mathbb{E}_{ x \sim \mathcal{D},\, \{y_i\}_{i=1}^G \sim \pi_{\theta_\text{old}}( \cdot | x) }
\left[ \frac{1}{G} \sum_{i=1}^{G} \frac{1}{|y_i|} \sum_{t=1}^{|y_i|} 
\min \left( w_{i,t}(\theta) \widehat{A}_{i,t},  \, \mathrm{clip} \left( w_{i,t}(\theta), 1 - {\varepsilon}, 1 + {\varepsilon}\right) \widehat{A}_{i,t} \right)
\right],
\end{align}

Here, $G$ denotes the total number of responses generated for every query $x$ (referred to as the group size), while the importance ratio $w_{i,t}(\theta)$ and the advantage estimate $\widehat{A}_{i,t}$ corresponding to token $y_{i,t}$ are defined as follows:

\begin{align}
    w_{i,t}(\theta)=\frac{ \pi_{\theta} (y_{i,t} | x, y_{i,<t}) }{ \pi_{\theta_\text{old}} (y_{i,t} | x,y_{i,<t})},\quad\ 
    \widehat{A}_{i,t} = \widehat{A}_{i} = \frac{r(x, y_i) - \mathrm{mean} \left( \{ r(x, y_i) \}_{i=1}^G \right) }{ \mathrm{std} \left( \{ r(x, y_i) \}_{i=1}^G \right) },
\end{align}

respectively, with each token in $y_i$ possessing the identical advantage value $\widehat{A}_{i}$.

\section{Motivation}
\label{sec:motivation}

As model architectures continue to expand in terms of size, incorporate increased sparsity (such as in Mixture-of-Experts models), and generate longer outputs, the requirement for a substantial rollout batch size becomes paramount in order to fully leverage hardware capabilities during reinforcement learning (RL). In order to enhance sample efficiency, it is customary to divide this large rollout batch into several smaller mini-batches, each of which is used for separate gradient update steps. However, this approach unavoidably leads to an off-policy training scenario, since the responses $y$ are drawn from a previous policy $\pi_{\theta_\text{old}}$, and not from the current policy $\pi_\theta$ undergoing optimization. Consequently, this highlights the importance of the clipping techniques in algorithms like PPO and GRPO, which serve to filter out samples that are excessively ``off-policy'' and thereby safeguard the stability of gradient computation.

Although techniques such as clipping are employed to address off-policy discrepancies, we observe a deeper, underlying flaw in GRPO: \textit{the objective itself is ill-posed}. This issue is especially pronounced during the training of large-scale models on tasks that require extended responses, often resulting in catastrophic collapse of the model. The root cause of the ill-posed GRPO objective lies in an improper use of importance sampling weights. Importance sampling is fundamentally intended to compute the expectation of a function $f$ under a target distribution $\pi_\text{tar}$ by appropriately reweighting samples obtained from a behavior distribution $\pi_\text{beh}$:

\begin{align}
\mathbb{E}_{z \sim \pi_\text{tar}} \left[ f(z) \right] = \mathbb{E}_{z \sim \pi_\text{beh}} \left[ \frac{ \pi_\text{tar} (z) }{ \pi_\text{beh} (z)} \, f(z) \right].
\end{align}

Importantly, effective correction for the distributional discrepancy using the importance weight $\frac{ \pi_\text{tar} (z) }{ \pi_\text{beh} (z)}$ requires averaging across a large number of samples ($N \gg 1$) drawn from the behavior distribution $\pi_\text{beh}$.

By comparison, GRPO assigns the importance weight $\frac{ \pi_{\theta} (y_{i,t} | x, y_{i,<t}) }{ \pi_{\theta_\text{old}} (y_{i,t} | x,y_{i,<t})}$ individually at every token step $t$. This weighting relies on a single sampled token $y_{i,t}$ from each corresponding next-token distribution $\pi_{\theta_\text{old}}( \cdot | x, y_{i, <t})$, which undermines its intended function of correcting for distribution shifts. Rather than achieving the desired correction, this practice injects significant high-variance noise into gradient computations, leading to greater instability especially for long output sequences, a problem intensified further by the presence of clipping. Through empirical study, we have found that such issues can result in a catastrophic breakdown of the model—an outcome from which recovery is rarely possible. Efforts to resume training post-collapse—whether by restoring previous checkpoints, extensively tuning hyperparameters like clipping thresholds, increasing the generation horizon, or altering the nature of RL queries—consistently fail to reverse the damage.

The observation outlined above reveals a central flaw in the design of GRPO. Specifically, the ineffectiveness of importance weighting at the token level highlights an essential guideline: \textbf{the granularity of the optimization target must align with that of the reward}. Given that the reward is assigned per sequence, conducting off-policy correction on a per-token basis is inherently flawed. This insight leads us to abandon the token-level objective, and instead consider applying importance weighting and optimization directly at the \textit{sequence level}.

\section{Algorithm}

\subsection{GSPO: Group Sequence Policy Optimization}

Although the use of token-level importance weights $\frac{ \pi_{\theta} (y_{i,t} | x, y_{i,<t}) }{ \pi_{\theta_\text{old}} (y_{i,t} | x,y_{i,<t})}$ introduces issues in GRPO, we find that, for language generation, the \textit{sequence-level} importance weight $\frac{ \pi_{\theta} (y | x) }{ \pi_{\theta_\text{old}} (y | x)}$ possesses a well-defined theoretical interpretation: it quantifies the discrepancy between the output $y$, which is drawn from $\pi_{\theta_\text{old}} (\cdot | x)$, and the distribution $\pi_{\theta} (\cdot | x)$. This correspondence is in direct agreement with sequence-level rewards and can effectively function as an informative metric for the application of the clipping strategy.

Building upon this simple insight, we introduce the \textbf{Group Sequence Policy Optimization (GSPO)} algorithm. GSPO is guided by the subsequent optimization objective defined at the sequence level:

\begin{align}
\mathcal{J}_\text{GSPO} (\theta) =
\mathbb{E}_{ x \sim \mathcal{D},\, \{y_i\}_{i=1}^G \sim \pi_{\theta_\text{old}}( \cdot | x) }
\left[ 
\frac{1}{G} \sum_{i=1}^{G}
\min \left( s_{i}(\theta)  \widehat{A}_{i},  \, \mathrm{clip} \left( s_{i}(\theta), 1 - {\varepsilon}, 1 + {\varepsilon} \right) \widehat{A}_{i} \right) 
\right],
\label{equ:gspo}
\end{align}

where we utilize advantage estimation based on groups:

\begin{align}
\widehat{A}_{i} = \frac{r(x, y_i) - \mathrm{mean} \left( \{ r(x, y_i) \}_{i=1}^G \right) }{ \mathrm{std} \left( \{ r(x, y_i) \}_{i=1}^G \right) },
\end{align}

and construct the importance ratio $s_{i}(\theta)$ utilizing sequence likelihood as described by \citep{click}:

\begin{align}
s_{i}(\theta) = \left( \frac{ \pi_{\theta} (y_i | x) }{ \pi_{\theta_\text{old}} (y_i | x)} \right)^{\frac{1}{|y_i|}}
=
\exp \left( \frac{1}{|y_i|} \sum_{t=1}^{|y_i|} \log \frac{ \pi_{\theta} (y_{i,t} | x, y_{i,<t}) }{ \pi_{\theta_\text{old}} (y_{i,t} | x,y_{i,<t})} \right).
\end{align}

Consequently, GSPO implements clipping at the full response level rather than at the level of individual tokens. This approach serves to filter out significantly "off-policy" samples during the estimation of gradients, thereby ensuring compatibility with both sequence-level reward mechanisms and optimization processes. Importantly, we incorporate length normalization in $s_{i}(\theta)$, which helps mitigate variance and confines $s_{i}(\theta)$ to a consistent numerical scale. Without normalization, large shifts in the likelihood of a handful of tokens can induce substantial volatility in the sequence-level importance ratio, making it necessary to use different clipping ranges for responses of varying lengths. Additionally, it should be emphasized that the magnitudes of the clipping ranges used in GSPO differ, often significantly, from those in earlier approaches (such as GRPO) because of fundamental differences in how importance ratios are defined.

\subsection{Gradient Analysis}
\label{subsec:gradient}

The gradient of the GSPO objective can be computed in the following manner (for clarity, clipping has been excluded):

\begin{align}
\nabla_{\theta} \mathcal{J}_\text{GSPO} (\theta)
=&\ 
\nabla_{\theta} \mathbb{E}_{ x \sim \mathcal{D},\, \{y_i\}_{i=1}^G \sim \pi_{\theta_\text{old}}( \cdot | x) }
\left[ 
\frac{1}{G} \sum_{i=1}^{G} s_{i}(\theta) \widehat{A}_{i}
\right] \\
= &\ 
\mathbb{E}_{ x \sim \mathcal{D},\, \{y_i\}_{i=1}^G \sim \pi_{\theta_\text{old}}( \cdot | x) }
\left[ 
\frac{1}{G} \sum_{i=1}^{G}
s_{i}(\theta) \widehat{A}_{i}
\cdot \nabla_{\theta} \log s_{i}(\theta)
\right] \\
=&\ 
\mathbb{E}_{ x \sim \mathcal{D},\, \{y_i\}_{i=1}^G \sim \pi_{\theta_\text{old}}( \cdot | x) }
\left[ 
\frac{1}{G} \sum_{i=1}^{G} 
\left( \frac{ \pi_{\theta} (y_i | x) }{ \pi_{\theta_\text{old}} (y_i | x)} \right)^{\frac{1}{|y_i|}} \widehat{A}_{i} 
\cdot \frac{1}{|y_i|} \sum_{t=1}^{|y_i|} \nabla_{\theta} \log \pi_{\theta} (y_{i,t} | x, y_{i,<t}) 
\right].
\label{equ:gspo_grad}
\end{align}

For reference, the gradient for the GRPO objective can be written as (observe that $\widehat{A}_{i,t} = \widehat{A}_{i}$):

\begin{align}
\nabla_{\theta} \mathcal{J}_\text{GRPO} (\theta) 
=&\ 
\nabla_{\theta} \mathbb{E}_{ x \sim \mathcal{D},\, \{y_i\}_{i=1}^G \sim \pi_{\theta_\text{old}}( \cdot | x) }
\left[ \frac{1}{G} \sum_{i=1}^{G} \frac{1}{|y_i|} \sum_{t=1}^{|y_i|} 
w_{i,t}(\theta) \widehat{A}_{i,t}
\right] \\
=&\
\mathbb{E}_{ x \sim \mathcal{D},\, \{y_i\}_{i=1}^G \sim \pi_{\theta_\text{old}}( \cdot | x) }
\left[ \frac{1}{G} \sum_{i=1}^{G} \widehat{A}_{i} 
\cdot \frac{1}{|y_i|} \sum_{t=1}^{|y_i|} 
\frac{ \pi_{\theta} (y_{i,t} | x, y_{i,<t}) }{ \pi_{\theta_\text{old}} (y_{i,t} | x,y_{i,<t})} 
\nabla_{\theta} \log \pi_{\theta} (y_{i,t} | x, y_{i,<t})  
\right].
\end{align}

Accordingly, the principal difference between GSPO and GRPO concerns \textit{the manner in which the gradients of token log likelihoods are weighted}. In GRPO, the contribution of each token is scaled by its corresponding ``importance weight'' $\frac{ \pi_{\theta} (y_{i,t} | x, y_{i,<t}) }{ \pi_{\theta_\text{old}} (y_{i,t} | x,y_{i,<t})}$. These weights are not uniform—they can fall within the ranges $(0, 1+\varepsilon]$ when $\widehat{A}_{i} > 0$ or $[1-\varepsilon, +\infty)$ if $\widehat{A}_{i} < 0$—and thus are significant and can aggregate over time, potentially causing erratic behavior as training continues. By comparison, GSPO assigns equal weight to every token within a response, thereby removing the instability introduced by the variable weighting in GRPO.

\subsection{GSPO-token: A Token-level Objective Variant}

In contexts such as multi-turn RL, there may be a need for more granular adjustment of advantage beyond the sequence level. Therefore, we propose a token-level variant of GSPO, termed \textbf{GSPO-token}, which facilitates customization of the advantage on a per-token basis:

\begin{align}
\mathcal{J}_\text{GSPO-token}(\theta)
=
\mathbb{E}_{ x \sim \mathcal{D},\, \{y_i\}_{i=1}^G \sim \pi_{\theta_\text{old}}( \cdot | x) }
\left[ 
\frac{1}{G} \sum_{i=1}^{G} \frac{1}{|y_i|} \sum_{t=1}^{|y_i|} 
\min \left( s_{i,t}(\theta) \widehat{A}_{i,t},  \, \mathrm{clip} \left( s_{i,t}(\theta), 1 - {\varepsilon}, 1 + {\varepsilon} \right) \widehat{A}_{i,t} \right) 
\right],
\label{equ:gspo-token}
\end{align}

in which

\begin{align}
s_{i,t}(\theta) 
= \mathrm{sg} \left[ s_{i}(\theta) \right]  \cdot \frac{ \pi_{\theta} (y_{i,t} | x, y_{i,<t}) }{ \mathrm{sg} \left[ \pi_{\theta} (y_{i,t} | x, y_{i,<t}) \right] },
\end{align}

and $\mathrm{sg}[\cdot]$ indicates extracting the numerical value while blocking the propagation of gradients, which is equivalent to the \texttt{detach} function in PyTorch. The gradient associated with GSPO-token can thus be expressed as:

\begin{align}
\nabla_{\theta} \mathcal{J}_\text{GSPO-token}(\theta)
=&\ 
\nabla_{\theta} \mathbb{E}_{ x \sim \mathcal{D},\, \{y_i\}_{i=1}^G \sim \pi_{\theta_\text{old}}( \cdot | x) }
\left[ 
\frac{1}{G} \sum_{i=1}^{G}
\frac{1}{|y_i|} \sum_{t=1}^{|y_i|} 
s_{i,t}(\theta) \widehat{A}_{i,t}
\right] \\
=&\ 
\mathbb{E}_{ x \sim \mathcal{D},\, \{y_i\}_{i=1}^G \sim \pi_{\theta_\text{old}}( \cdot | x) }
\left[ 
\frac{1}{G} \sum_{i=1}^{G} s_i (\theta) \cdot \frac{1}{|y_i|} \sum_{t=1}^{|y_i|} 
\widehat{A}_{i,t} \frac{ \nabla_{\theta} \pi_{\theta} (y_{i,t} | x, y_{i,<t}) }{ \pi_{\theta} (y_{i,t} | x, y_{i,<t}) }
\right] \\
=&\ 
\mathbb{E}_{ x \sim \mathcal{D},\, \{y_i\}_{i=1}^G \sim \pi_{\theta_\text{old}}( \cdot | x) }
\left[ 
\frac{1}{G} \sum_{i=1}^{G}  \left( \frac{ \pi_{\theta} (y_i | x) }{ \pi_{\theta_\text{old}} (y_i | x)} \right)^{\frac{1}{|y_i|}}
\cdot \frac{1}{|y_i|} \sum_{t=1}^{|y_i|}
\widehat{A}_{i,t} \nabla_{\theta} \log \pi_{\theta} (y_{i,t} | x, y_{i,<t})  
\right].
\label{equ:gspo-token_grad}
\end{align}

Observe that the ratio $\frac{ \pi_{\theta} (y_{i,t} | x, y_{i,<t}) }{ \mathrm{sg} \left[ \pi_{\theta} (y_{i,t} | x, y_{i,<t}) \right] }$ evaluates to 1, meaning that $s_{i,t}(\theta)$ and $s_{i}(\theta)$ yield the same numerical result. If we compare Equation~\eqref{equ:gspo} with~\eqref{equ:gspo-token}, and Equation~\eqref{equ:gspo_grad} with~\eqref{equ:gspo-token_grad}, it becomes evident that GSPO-token and GSPO produce equivalent numerical results for the optimization target, clipping criterion, and analytical gradients under the condition that every token in the output $y_i$ shares an identical advantage (that is, $\widehat{A}_{i,t} = \widehat{A}_{i}$ for all $t$). However, GSPO-token offers additional versatility by permitting the assignment of distinct advantage values to each token.

\section{Experiments and Discussion}

\subsection{Empirical Results}

We conduct experiments using a cold-start model, fine-tuned from Qwen3-30B-A3B-Base, and present both training reward trajectories and performance metrics across the AIME'24 (average Pass@1 computed over 32 samples), LiveCodeBench (202410-202502, average Pass@1 computed over 8 samples), and CodeForces (Elo Rating) evaluation suites. Throughout the RL training procedure, each batch of rollout data is subdivided into four mini-batches to facilitate gradient updates. For the GSPO approach, the clipping boundaries in Equation~\eqref{equ:gspo} are set to 3e-4 (left) and 4e-4 (right). As a point of comparison, we use GRPO as the baseline, tuning the left and right clipping intervals in Equation~\eqref{equ:grpo} to 0.2 and 0.27, respectively; these values were selected with careful calibration to guarantee a fair comparison. It is important to highlight that GRPO requires employing the Routing Replay training technique for MoE RL to converge reliably, a topic explored further in \S~\ref{subsec:routing_replay}, whereas \textit{GSPO eliminates the dependence on this approach}.

As depicted in Figure~\ref{fig:results}, the training process using GSPO remains consistently stable. Notably, \textbf{GSPO facilitates ongoing enhancement in performance by scaling up the training compute, periodically refreshing the query set, and increasing the generation length}. Additionally, GSPO surpasses GRPO in terms of training efficiency, achieving higher accuracy and superior benchmark outcomes while utilizing equivalent training compute and an equal number of queries. Lastly, our successful implementation of GSPO in the RL training of the most recent Qwen3 models provides robust evidence for the effectiveness of GSPO in leveraging RL scaling capabilities for large language models.

\subsection{Curious Observation on Clipping Fractions}

An important differentiation between GSPO and GRPO lies in their approach to response clipping: GSPO applies clipping at the sequence level, whereas GRPO does so at the individual token level. Notably, as illustrated in Figure~\ref{fig:clipping}, there exists a gap of two orders of magnitude in the proportion of tokens subjected to clipping by GSPO compared to GRPO (this substantial difference persists even when the clipping ranges are modified). Nonetheless, despite GSPO discarding a far greater number of tokens and thereby leveraging a smaller set for training or gradient computation, it consistently demonstrates higher training efficiency relative to GRPO. This somewhat surprising outcome — where greater token clipping yields enhanced training efficiency — implies that the token-level gradient estimates in GRPO are prone to noise and lack effectiveness for exploiting samples. By contrast, GSPO’s reliance on sequence-level gradients offers a more robust and informative learning signal.

\subsection{Benefit of GSPO for MoE Training}
\label{subsec:routing_replay}

\paragraph{Background}
Relative to the RL training process in dense models, MoE models face distinctive stability problems due to their sparse activation patterns. Specifically, our observations indicate that when the GRPO algorithm is employed, MoE models can exhibit significant \textit{expert-activation volatility}, which obstructs the convergence of RL training. To elaborate, the subset of experts engaged in producing the same output may shift substantially after performing one or more gradient steps. For instance, considering the 48-layer Qwen3-30B-A3B-Base model, it was observed that, for a given rollout sample, about 10\% of the experts triggered post-RL gradient update (under the revised policy $\pi_{\theta}$) differ from those selected previously (under the original policy $\pi_{\theta_\text{old}}$). This instability, which intensifies with deeper MoE architectures, leads to pronounced oscillations in the token-level importance weights $w_{i,t}(\theta) = \frac{ \pi_{\theta} (y_{i,t} | x, y_{i,<t}) }{ \pi_{\theta_\text{old}} (y_{i,t} | x,y_{i,<t})}$, as described in \S~\ref{sec:motivation} and~\ref{subsec:gradient}. Such drastic fluctuations undermine the reliability of these weights and, as a result, disrupt the expected convergence behavior during RL training.

\paragraph{Our Previous Approach} In our earlier work, we addressed this issue by utilizing the \textbf{Routing Replay} method during training. This involved storing the expert routing decisions generated by $\pi_{\theta_\text{old}}$, then reusing these routing selections within $\pi_{\theta}$ as we computed the importance ratios $w_{i,t}(\theta) = \frac{ \pi_{\theta} (y_{i,t} | x, y_{i,<t}) }{ \pi_{\theta_\text{old}} (y_{i,t} | x,y_{i,<t})}$. This approach guarantees that, for any given token $y_{i,t}$, both $\pi_{\theta} (y_{i,t} | x, y_{i,<t})$ and $\pi_{\theta_\text{old}} (y_{i,t} | x,y_{i,<t})$ utilize an identical set of activated experts. Consequently, the token-level importance ratios become more stable, and it is possible to optimize the same expert network consistently across gradient updates. As shown in Figure~\ref{fig:routing_replay}, Routing Replay is critical for enabling stable and effective convergence of GRPO training in MoE models.

\paragraph{Benefit of GSPO} While Routing Replay permits the effective GRPO training of MoE models, its approach of retaining routing strategies introduces supplementary memory consumption, increases communication demands, and may constrain the attainable capacity of the MoE model. By contrast, as illustrated in Figure~\ref{fig:results}, GSPO dispenses with Routing Replay entirely, operating without such dependencies. GSPO can conventionally compute the importance ratios $s_i(\theta)$, thereby ensuring both convergence and training stability. Crucially, GSPO restricts its focus to sequence-level likelihoods (i.e., $\pi_{\theta} (y_i | x)$) and remains unaffected by fluctuations in token-level likelihoods (i.e., $\pi_{\theta} (y_{i,t} | x, y_{i,<t})$). As the MoE model consistently upholds its language modeling proficiency, variations in sequence likelihood are kept minimal. Consequently, GSPO directly addresses the issue of instability in expert activation within MoE models, eliminating the necessity for complicated methods such as Routing Replay. This direct approach not only streamlines and fortifies training, but also enables the exploitation of the model’s full representational capacity without imposed limitations.

\subsection{Benefit of GSPO for RL Infrastructure}

Due to the precision mismatches that exist between training engines such as Megatron and inference engines like SGLang and vLLM, the standard approach is to recompute the likelihoods of generated responses under the previous policy $\pi_{\theta_\text{old}}$ using the training engine. In contrast, GSPO relies on sequence-level likelihoods during optimization, rather than on token-level likelihoods. Sequence-level measures are generally far less sensitive to precision differences between engines. As a result, GSPO enables direct utilization of the likelihood values produced by the inference engine for optimization purposes, eliminating the need for recomputation with the training engine. This feature is particularly advantageous for use cases such as partial rollout, multi-turn RL, and for systems where training and inference are decoupled.

\section{Conclusion}

We introduce Group Sequence Policy Optimization (GSPO), a novel reinforcement learning framework designed for optimizing large language models. GSPO builds on the concept of importance sampling by formulating importance ratios grounded in sequence likelihood, and applies sequence-level mechanisms for clipping, reward assignment, and policy optimization. Relative to GRPO, GSPO offers marked enhancements in training robustness, computational efficiency, and overall effectiveness, and is especially well-suited for extensive RL training of Mixture of Experts (MoE) architectures. This establishes GSPO as the key driver behind the notable advancements observed in the recent Qwen3 models. As GSPO provides a scalable methodological foundation, we aim to further expand reinforcement learning capacities and anticipate significant progress in the development of advanced intelligence.

\end{document}